\subsection{Workload Model and Online Duration Estimation}
\label{sec:model}
A Draft Sequence's duration depends on its stages and how fast they execute under current conditions.
The controller therefore learns stage durations from completed Draft Sequences and combines them according to the sequence being considered.

\smallskip\noindent\textbf{Modeled workload sequence.}
Each combination of a processing stage and its configuration is identified by a workload key \(k\).
Capture \(i\)'s Draft Sequence is represented by the ordered keys \(\mathcal{K}_i=(k_1,\ldots,k_m)\).
At position \(j\), the remaining sequence is \(\mathcal{K}_{i,j}=(k_j,\ldots,k_m)\).
The last key covers mandatory processing from encoding through Draft image saving.
Including this processing lets admission account for the time needed to produce the image after the optional stages finish.

\smallskip\noindent\textbf{Online point estimation.}
The controller maintains a baseline duration \(\hat b(k)\) for each key to represent its long-term processing cost.
A shared factor \(\gamma\) adjusts these baselines to recent execution conditions.
Multiplying the two gives the predicted duration \(\hat p(k)\), and summing these predictions gives the duration estimate for a modeled sequence \(\mathcal K\):
\begin{equation}
  \hat P(\mathcal{K})
  =\sum_{k\in\mathcal{K}}\hat p(k),
  \qquad
  \hat p(k)=\hat b(k)\,\gamma .
  \label{eq:point-estimate}
\end{equation}
At Draft Sequence completion, the controller divides each observed stage duration by the shared factor available before the update to account for the previously estimated execution conditions.
It incorporates this adjusted duration into the key's cumulative average to update \(\hat b(k)\).

To update \(\gamma\), the controller divides each observed duration by the corresponding baseline available before the update and takes the median of these ratios.
Only keys with an existing baseline contribute; if none are available, \(\gamma\) remains unchanged.
The median limits the influence of isolated stage slowdowns, while the shared factor carries recent changes in execution speed into subsequent estimates.

An unobserved key has a zero estimate until its first duration is recorded at Draft Sequence completion.
Section~\ref{sec:admission} explains how admission obtains these initial observations.
The learned baselines and shared factor are retained when the queue drains, so subsequent captures can reuse the estimates.
